% !TEX root = ../../1Main.tex
% Aufgabenstellung - Template

% ===================================================================
% ANLEITUNG: Dieses Dokument ist die offizielle Aufgabenstellung
% Ihrer Arbeit. Normalerweise wird dies von Ihrem Betreuer erstellt.
% Aktivieren Sie dies in 1Main.tex nur wenn Sie eine formale
% Aufgabenstellung haben!
% ===================================================================

% === Sauberer Seitenumbruch ===
\clearpage
\cleardoublepage

% === Spezielle Formatierung für Aufgabenstellung ===
\begingroup
\let\cleardoublepage\relax
\let\clearpage\relax

% === Institution Logo ===
\begin{figure}[t]
  \vspace{1.5cm}
  \begin{flushright}
    \IfFileExists{default_style/AtheneGespiegelt.png}{%
      \includegraphics[width=0.1\textwidth]{default_style/AtheneGespiegelt.png}%
    }{%
      \textcolor{secondary}{\rule{0.1\textwidth}{0.1\textwidth}}%
    }
  \end{flushright}
  \vspace{-4.5cm}
\end{figure}

% === Aufgabenstellung ===
\chapter*{Aufgabenstellung}

\section*{Hintergrund}
\textit{[Beschreiben Sie hier den fachlichen und wissenschaftlichen Hintergrund Ihrer Arbeit. Warum ist das Thema relevant? Welche Probleme sollen gelöst werden?]}

\textit{[Kontext und Motivation der Arbeit...]}

\newpage

\section*{Aufgaben}

\begin{enumerate}
    \item \textbf{Motivation und Problemstellung:}\\
      \textit{[Beschreibung der Motivation, konkrete Problemstellung, wissenschaftliche Fragestellungen, Begriffsdefinitionen]}
    
    \item \textbf{Grundlagenanalyse und Literaturrecherche:}\\
      \textit{[Umfang der Literaturrecherche, zu analysierende Quellen, theoretische Einordnung]}

    \item \textbf{Konzeption und Methodik:}\\
      \textit{[Methodischer Ansatz, Architektur, Designentscheidungen]}

    \item \textbf{Implementierung/Durchführung:}\\
      \textit{[Konkrete Umsetzungsaufgaben, technische Anforderungen, zu entwickelnde Komponenten]}

    \item \textbf{Evaluation:}\\
      \textit{[Bewertungskriterien, Tests, Validierung der Ergebnisse]}

    \item \textbf{Zusammenfassung und Ausblick:}\\
      \textit{[Zusammenfassung der Erkenntnisse, kritische Reflexion, zukünftige Arbeiten]}
  \end{enumerate}

\endgroup

% === Seitenfüllung und Abschluss ===
\vspace*{\fill}
\clearpage
\thispagestyle{empty}

% ===================================================================
% HINWEIS: Dieses Dokument ist optional! Löschen Sie es wenn Sie
% keine formale Aufgabenstellung benötigen.
% ===================================================================